% lstar_problem.tex -- a fully self-contained statement of the
% inequality L* for external mathematicians (approximation theory,
% orthogonal polynomials, sampling theory, analytic number theory,
% operator theory).  No project-internal vocabulary.  Builds with
% pdflatex.  Companion machine check: verify_lstar_instance.py.
\documentclass[11pt]{article}

\usepackage[a4paper,margin=2.7cm]{geometry}
\usepackage{amsmath,amssymb,amsthm}
\usepackage{booktabs}
\usepackage{microtype}
\usepackage{xcolor}
\usepackage[colorlinks=true,linkcolor=blue!50!black,urlcolor=blue!50!black]{hyperref}

\newtheorem{problem}{Problem}
\newtheorem{prop}{Proposition}
\theoremstyle{remark}
\newtheorem{remark}{Remark}

\newcommand{\dd}{\,\mathrm{d}}
\newcommand{\R}{\mathbb{R}}
\newcommand{\N}{\mathbb{N}}
\newcommand{\E}{\mathsf{E}}

\title{\bfseries A strict polynomial subordination inequality for two
explicit prime-generated discrete measures\\[2mm]
\large An open problem, stated self-contained}
\author{Stefan Hamann}
\date{August 25, 2026}

\begin{document}
\maketitle

\begin{abstract}\noindent
We define, completely explicitly, a family of pairs $(\mu,\nu)$ of
finite positive discrete measures on $[-1,1]$, built from the von
Mangoldt function through a fixed elementary sampling recipe.  Writing
$S$ for the total number of atoms of the pair, the problem is to prove
or refute:
\[
\int p^{2}\dd\nu \;<\; \int p^{2}\dd\mu
\qquad\text{for every real polynomial } p\neq 0 \text{ with }
\deg p < \tfrac{S+1}{2}.
\]
Numerically the inequality holds on all $42$ instances of the family
that we have tested, with a spectral margin that shrinks from
$1.7\cdot 10^{-4}$ to $1.4\cdot 10^{-7}$ as $S$ grows from $367$ to
$1755$, and it genuinely fails a few degrees beyond the stated bound.
Structurally destroyed comparison inputs lose the inequality already
at $11$--$15\%$ of the critical degree.  The degree bound
$(S+1)/2$ is exactly the moment-freedom count of an $S$-atom signed
measure.  No proof is known; a counterexample at large $S$ would be
equally valuable.  This document is a problem statement, not a claim.
\end{abstract}

\medskip\noindent
\textbf{Status note (August 27, 2026).}\enspace The research programme
that produced this problem has paused its internal work on it and
frozen it as a specialist problem.  The new
Section~\ref{sec:frozen} carries the current state: the measurement
campaign since the update of \S\ref{sec:numerics} has compressed the
problem into \emph{one object} --- an almost perfect joint wander of
two fine structures whose difference is five hundred times smaller
than either --- and three precise questions.
Sections~\ref{sec:problem}--\ref{sec:context} are unchanged.

\section{The problem}\label{sec:problem}

\subsection{The construction}\label{sec:construction}

Let $\Lambda$ denote the von Mangoldt function
($\Lambda(n)=\log q$ if $n=q^{m}$ is a prime power, else $0$), and let
$\gamma$ be Euler's constant.  Fix once and for all the
\emph{oversampling parameter} $\nu_{0}=4$ and the table cap
$T = 4\cdot 10^{5}$ (all prime powers used below lie far below $T$;
the cap only delimits the admissible index list in
\S\ref{sec:family}).

Fix a prime power $z$ (the \emph{window anchor}) and let $z'$ be the
next larger prime power.  Define
\[
\alpha=\log z,\qquad
\delta=\frac{\log z'-\log z}{2\nu_{0}},\qquad
M=\Big\lceil \frac{\alpha}{\delta}-10^{-9}\Big\rceil+1
\ \ (\text{increased by }1\text{ if odd}),
\]
\[
\Delta=\frac{2\alpha}{M},\qquad
L=2M-2,\qquad
N=\frac{M}{2}.
\]
Let $t_{\Delta}(x)=\max(0,\,1-|x|/\Delta)$ be the unit triangle
(``tent'') of half-width $\Delta$.

\medskip\noindent
\textbf{Step 1 (lag sequence).}  For $i=0,1,\dots,M-1$ set
$c_{i}=A(i\Delta)+P(i\Delta)$, where
\begin{align}
P(s) &= -\sum_{n\le z^{2}}\frac{\Lambda(n)}{\sqrt{n}}
  \Bigl[t_{\Delta}(s-\log n)+t_{\Delta}(s+\log n)\Bigr],
  \label{eq:P}\\[1mm]
A(s) &= -(\gamma+\log\pi)\,t_{\Delta}(s)
  +\int_{0}^{\infty}
  \frac{2\,t_{\Delta}(s)\,e^{-2w}
        -\bigl[t_{\Delta}(s-w)+t_{\Delta}(s+w)\bigr]e^{-w/2}}
       {1-e^{-2w}}\;\dd w .
  \label{eq:A}
\end{align}
Both integrands are piecewise smooth and absolutely integrable ($A$ is
the archimedean term of the Weil explicit formula evaluated on the
tent test function $t_{\Delta}(\cdot-s)$; $P$ is the triangular-window
sampling of the prime comb --- but the two formulas above are the
complete definition, no outside input is needed).

\medskip\noindent
\textbf{Step 2 (spectral density of the symmetrized lags).}  Define
the even trigonometric polynomial
\begin{equation}\label{eq:symbol}
f(\theta)\;=\;c_{0}+2\sum_{k=1}^{M-2}c_{k}\cos(k\theta)
              +c_{M-1}\cos\bigl((M-1)\theta\bigr),
\end{equation}
i.e.\ the discrete Fourier transform of the symmetric extension
$(c_{0},\dots,c_{M-1},c_{M-2},\dots,c_{1})$ of the lag sequence; the
values $d_{j}=f(\theta_{j})$ at $\theta_{j}=2\pi j/L$ are the
eigenvalues of the circulant extension of the $L\times L$ symmetric
Toeplitz matrix with symbol~$f$.

\medskip\noindent
\textbf{Step 3 (folding to $[-1,1]$ and sign split).}
For $j=1,\dots,L/2$ put
\begin{equation}\label{eq:atoms}
x_{j}=\cos\theta_{j},\qquad
\widetilde w_{j}
  =\kappa_{j}\,\frac{2}{L}\,\bigl(1-\cos\theta_{j}\bigr)\,
   f(\theta_{j}),
\qquad
\kappa_{j}=\begin{cases}1,& 1\le j<L/2,\\[0.5mm]
\tfrac12,& j=L/2 .\end{cases}
\end{equation}
Discard indices with $\widetilde w_{j}=0$ and let $S$ be the number of
remaining atoms.  Define the two positive measures
\begin{equation}\label{eq:munu}
\mu \;=\; \sum_{\widetilde w_{j}>0}\widetilde w_{j}\,\delta_{x_{j}},
\qquad
\nu \;=\; \sum_{\widetilde w_{j}<0}\bigl(-\widetilde w_{j}\bigr)\,
          \delta_{x_{j}},
\end{equation}
with atom counts $S_{+}$ and $S_{-}$, $S_{+}+S_{-}=S$.  Both are
finite positive measures with finite (disjoint) supports in $[-1,1)$.

\subsection{The concrete running example \texorpdfstring{($z=16$)}{(z=16)}}
\label{sec:example}

For $z=16$ (so $z'=17$, $\alpha=\log 16=2.772589\ldots$,
$\delta=\tfrac18\log\tfrac{17}{16}=0.00757808\ldots$) the recipe gives
\[
M=368,\quad \Delta=0.01506842\ldots,\quad L=734,\quad N=184,
\]
$70$ prime powers $n\le 256$ enter \eqref{eq:P}, and the fold produces
\[
S=367,\qquad S_{+}=263,\qquad S_{-}=104,\qquad
\mu(\R)=3.699410\ldots,\qquad \nu(\R)=0.226205\ldots
\]
(every one of the $367$ candidate atoms carries a nonzero weight, so
$N=(S+1)/2$).  Table~\ref{tab:atoms} shows an excerpt; the companion
script (\S\ref{sec:script}) prints every atom.

\begin{table}[ht]
\centering
\caption{Atom excerpt for $z=16$ ($L=734$; $x_{j}=\cos(2\pi j/L)$).
Positive weights belong to $\mu$, negative weights (in absolute value)
to $\nu$.}
\label{tab:atoms}
\begin{tabular}{rrrrl}
\toprule
$j$ & $x_{j}$ & $f(\theta_{j})$ & $\widetilde w_{j}$ & channel\\
\midrule
$1$   & $+0.999963$ & $+27.427975$ & $+2.738\cdot10^{-6}$ & $\mu$\\
$2$   & $+0.999853$ & $-10.075060$ & $-4.023\cdot10^{-6}$ & $\nu$\\
$3$   & $+0.999670$ & $+4.824158$  & $+4.334\cdot10^{-6}$ & $\mu$\\
$4$   & $+0.999414$ & $-2.729923$  & $-4.360\cdot10^{-6}$ & $\nu$\\
$5$   & $+0.999084$ & $+1.709341$  & $+4.266\cdot10^{-6}$ & $\mu$\\
$6$   & $+0.998681$ & $-1.136329$  & $-4.083\cdot10^{-6}$ & $\nu$\\
\multicolumn{5}{c}{$\cdots$}\\
$183$ & $+0.004280$ & $+8.446382$  & $+2.292\cdot10^{-2}$ & $\mu$\\
$184$ & $-0.004280$ & $-2.274690$  & $-6.225\cdot10^{-3}$ & $\nu$\\
$185$ & $-0.012840$ & $+6.118735$  & $+1.689\cdot10^{-2}$ & $\mu$\\
\multicolumn{5}{c}{$\cdots$}\\
$365$ & $-0.999853$ & $+4.611240$  & $+2.513\cdot10^{-2}$ & $\mu$\\
$366$ & $-0.999963$ & $+6.747708$  & $+3.677\cdot10^{-2}$ & $\mu$\\
$367$ & $-1.000000$ & $-1.223818$  & $-3.335\cdot10^{-3}$ & $\nu$\\
\bottomrule
\end{tabular}
\end{table}

\subsection{The admissible family and the question}\label{sec:family}

Call the anchor $z$ \emph{admissible} if $N=M/2\in[128,900]$ and at
least $40$ prime powers lie below $z^{2}$ (with $z$ drawn from the
prime powers below $3.8\cdot10^{5}$, well inside the cap $T$).
Exactly $42$ prime powers are admissible,
\begin{align*}
z\in\{&16,\,23,\,25,\,27,\,29,\,31,\,37,\,41,\,43,\,47,\,49,\,53,\,
59,\,61,\,64,\,67,\,71,\,73,\\
&79,\,81,\,83,\,89,\,97,\,103,\,109,\,113,\,121,\,128,\,131,\,139,\,
151,\,157,\,163,\,169,\\
&173,\,181,\,199,\,211,\,233,\,243,\,293,\,317\},
\end{align*}
with $S$ ranging over $[283,1755]$ and $N$ over $[142,878]$.  We call
$\{(\mu_z,\nu_z)\}$ the \emph{prime window family} and $N=(S+1)/2$ the
\emph{critical degree} of an instance.

\begin{problem}[half-filling subordination]\label{prob:main}
Prove or refute: for every admissible anchor $z$ and every real
polynomial $p\neq 0$ with $\deg p<\frac{S+1}{2}$,
\begin{equation}\label{eq:lstar}
\int p^{2}\dd\nu_{z}\;<\;\int p^{2}\dd\mu_{z}.
\tag{$L^{*}$}
\end{equation}
\end{problem}

Any of the following would be decisive progress: (a) a proof of
\eqref{eq:lstar} for the full family, or for the unbounded family with
the cap $N\le 900$ removed; (b) a structural proof for a single
instance (say $z=16$) that explains the survival \emph{exactly up to}
the critical degree by a mechanism rather than by numerical
evaluation; (c) a counterexample --- an admissible $z$ (with or
without the cap) and an explicit $p$ violating \eqref{eq:lstar}.
Since each instance is a statement about finitely many explicitly
computable real numbers, a certified interval-arithmetic evaluation
also settles any single instance rigorously; the interest is in the
mechanism that makes the bound $(S+1)/2$ sharp, uniformly in the
family.

\section{Equivalent formulations}\label{sec:equiv}

Throughout fix an instance and write $\mu,\nu,S,N$ as above, and
$y_{1},\dots,y_{S_{-}}$ with weights $v_{k}>0$ for the atoms of
$\nu$.

\paragraph{(a) Hankel positive definiteness.}
Let $m_{k}=\int x^{k}\dd(\mu-\nu)$ and
$H_{n}=(m_{i+j})_{i,j=0}^{n-1}$.  For $p(x)=\sum_{i<n}a_{i}x^{i}$ one
has $\int p^{2}\dd(\mu-\nu)=a^{\mathsf T}H_{n}a$, so
\eqref{eq:lstar} is exactly
\[
H_{N}\succ 0
\qquad\text{with } N=\tfrac{S+1}{2},
\]
positive definiteness of the half-filling Hankel matrix of the
\emph{signed} measure $\mu-\nu$.  Equivalently (Jacobi), all leading
principal minors $D_{1},\dots,D_{N}$ of $H_{N}$ are positive, i.e.\
the three-term recurrence of $\mu-\nu$ has positive ``pivots''
$h_{0},\dots,h_{N-1}$ (quasi-definiteness through the window).

\paragraph{(b) Christoffel--Darboux kernel as a strict contraction.}
Let $(p_{i})_{i\ge 0}$ be the orthonormal polynomials of the
\emph{positive} measure $\mu$ and
$K_{n}(x,y)=\sum_{i<n}p_{i}(x)p_{i}(y)$ its reproducing
(Christoffel--Darboux) kernel.  Define the $S_{-}\times S_{-}$
symmetric matrix
\[
\E_{n}[k,l]\;=\;\sqrt{v_{k}v_{l}}\;K_{n}(y_{k},y_{l}),
\qquad 1\le k,l\le S_{-}.
\]
Writing $p=\sum_{i<n}a_{i}p_{i}$ and $B_{ki}=\sqrt{v_{k}}\,
p_{i}(y_{k})$ one has $\int p^{2}\dd\mu=\lVert a\rVert^{2}$ and
$\int p^{2}\dd\nu=\lVert Ba\rVert^{2}$, and
$\E_{n}=BB^{\mathsf T}$ has the same nonzero spectrum as
$B^{\mathsf T}B$.  Hence
\[
\eqref{eq:lstar}
\iff
\lambda_{\max}\bigl(\E_{N}\bigr)<1 :
\]
the $\nu$-dressed Christoffel--Darboux kernel of $\mu$ at the
half-filling degree is a strict contraction.  (Equivalently: the
Fredholm determinant $\det(I-s\,\E_{N})$ has no zero on $s\in(0,1]$.)

\paragraph{(c) Quadrature/sampling form.}
$\E_{n}$ is the Gram matrix, in $L^{2}(\mu)$, of the functions
$\sqrt{v_{k}}\,K_{n}(y_{k},\cdot)$; the problem asks whether the
sampling functionals $p\mapsto\sqrt{v_{k}}\,p(y_{k})$ at the $\nu$
atoms, viewed inside $L^{2}(\mu)$-normalized polynomials of degree
$<N$, form a strict Bessel system with bound $1$.  This places the
problem in the language of Marcinkiewicz--Zygmund/frame inequalities
for weighted node sets.

\section{What is known numerically}\label{sec:numerics}

All statements in this section are measurements (double precision,
with high-precision arbitration of every sign decision near zero);
none of them is a theorem.

\begin{enumerate}
\item \textbf{The inequality holds on all $42$ instances.}  The
margins $1-\lambda_{\max}(\E_{N})$ are all strictly positive and
\emph{decrease} with the instance size:
\begin{center}
\begin{tabular}{lrrr}
\toprule
$z$ & $S$ & $N$ & $1-\lambda_{\max}(\E_{N})$\\
\midrule
$16$  & $367$  & $184$ & $1.68\cdot10^{-4}$\\
$23$  & $301$  & $151$ & $7.65\cdot10^{-5}$\\
$31$  & $867$  & $434$ & $1.32\cdot10^{-5}$\\
$59$  & $979$  & $490$ & $2.09\cdot10^{-6}$\\
$128$ & $1677$ & $839$ & $1.82\cdot10^{-7}$\\
$233$ & $1717$ & $859$ & $1.42\cdot10^{-7}$ \quad(family minimum)\\
$317$ & $1067$ & $534$ & $1.72\cdot10^{-7}$\\
\bottomrule
\end{tabular}
\end{center}
The smallest instance ($z=16$) has the \emph{largest} margin; over the
tested range the margin loses roughly three orders of magnitude.  Any
asymptotic reading of the family should take this shrinkage seriously:
the unbounded version of Problem~\ref{prob:main} may fail at large
$S$, and a counterexample would be as informative as a proof.
\item \textbf{At the critical degree the contraction is broadly
tight.}  For $z=16$ the top of the spectrum of $\E_{184}$ is
\[
0.99983248,\quad 0.99874,\quad 0.99597,
\]
three eigenvalues above $0.995$ (confirmed by an independent
$60$-digit recomputation): the near-saturation is carried by several
directions, not by one isolated mode.
\item \textbf{The bound is sharp within $O(1)$ degrees.}  On every
instance, $\lambda_{\max}(\E_{n})$ passes $1$ at most six degrees
above $N$; the offset of the first Hankel-pivot sign failure beyond
$N$ has the census distribution
$\{0\!:\!18,\ 1\!:\!10,\ 2\!:\!6,\ 3\!:\!6,\ 4\!:\!1,\ 5\!:\!1\}$
over the $42$ instances.  For $z=16$ the failure is immediate:
$\lambda_{\max}(\E_{185})=1.00003660>1$.
\item \textbf{Structurally destroyed inputs fail early.}  Feeding the
same pipeline (Steps 1--3 verbatim) with destroyed input data ---
(i) the prime-comb positions replaced by i.i.d.\ uniform points on
$[0,2\alpha]$ at the same masses, (ii) the von Mangoldt weights
replaced by the Dirichlet-series coefficients of the Epstein zeta
function of $x^{2}+5y^{2}$, (iii) the comb replaced by its smooth
density, and (iv) a random pair-correlation model --- produces pairs
$(\mu,\nu)$ that lose \eqref{eq:lstar} already at degrees
$21/25/27/25$ respectively (at $N=184$), i.e.\ at $11$--$15\%$ of the
critical degree.  Whatever mechanism drives \eqref{eq:lstar} through
the full window, it is a property of the actual prime comb, not of
the sampling recipe.
\item \textbf{Norm bounds are far too coarse.}  On the $z=16$
instance the Gershgorin/row-sum bound on $\E_{n}$ already exceeds
$1$ at $n=21$, and the trace bound at $n=10$; both fail more than
$160$ degrees short of the truth.  A proof of \eqref{eq:lstar} must
exploit
cancellation, not absolute values.
\end{enumerate}

\subsection*{What is known (update, 2026-08-25)}

Three measured facts sharpen the picture above.

\emph{(a) The margin decay obeys a flattening power law and the family
was extended without a counterexample.}  The admissible family was
extended by $15$ further instances with critical degrees
$N=942..1218$ (about $40\%$ beyond the largest instance in the table);
\emph{every} margin $1-\lambda_{\max}(\E_{N})$ remains strictly
positive (minimum $+1.8\cdot10^{-8}$), each sign decision arbitrated
at $30$ and $45$ digits of working precision, so the inequality now
holds on all $57$ tested instances.  The decay follows a power law
$\mathrm{margin}\sim S^{-\alpha}$ with $\alpha\approx3.05$ that
\emph{flattens} at large $S$ (split-sample exponents $3.67/2.07$), and
the first-failure offset beyond $N$ stays $O(1)$ on the new instances
(maximum $+4$).  The margin also falls \emph{slower} than the local
per-degree growth rate of $\lambda_{\max}$ at the crossing --- the
decay is consistent with a fixed $O(1)$ degree offset at increasing
eigenvalue density rather than with an approaching failure.  The
unbounded version of Problem~\ref{prob:main} remains open, but the
tested range shows no crossing tendency.

\emph{(b) The problem is not diophantine.}  One might suspect that
\eqref{eq:lstar} consumes the number-theoretic (linear-forms-in-
logarithms) structure of the node positions $\log n$.  It does not:
replacing \emph{every} position $\log n$ by a nearby rational multiple
of the grid step (continued-fraction convergents, denominators up to
$56801$, position change at most $2.1\cdot10^{-9}$, weights and grid
cells preserved bitwise) --- which destroys every exact linear-form
relation among the positions --- reproduces every measured quantity of
the $z=16$ instance \emph{identically} ($\lambda_{\max}(\E_{184})=
0.99983248<1<1.00003660=\lambda_{\max}(\E_{185})$, confirmed at
$60$ digits).  Baker--Matveev-type nonresonance input is therefore not
what the inequality runs on (a needs-analysis puts the required
exponent about four orders of magnitude beyond what such bounds
deliver --- and the rational surrogate shows it is not needed at all
on the tested scales).

\emph{(c) The mechanism is metrically sharp.}  What the inequality
\emph{does} depend on is the exact metric placement of the nodes at
sub-grid precision: randomizing the fractional parts of the node
positions within their grid cells (distribution preserved exactly)
destroys \eqref{eq:lstar} once the perturbation exceeds roughly
$10^{-3}..3\cdot10^{-3}$ of the local node gap, and a plain jitter
ladder predicts the loss quantitatively.  Any proof therefore has to
consume the fractional-part \emph{profile} of the positions as a
metric object at that precision --- but, by (b), nothing finer.

\section{Why the degree bound matters}\label{sec:degree}

The critical degree is exactly the \emph{moment-freedom count}.  An
$S$-atom signed measure has exactly $S$ free moments
$m_{0},\dots,m_{S-1}$: the map (weights) $\mapsto$ (moment prefix) is
a Vandermonde bijection, and every $m_{k}$ with $k\ge S$ is forced
linearly through the recurrence whose coefficients are the node
polynomial $\prod_{j}(x-x_{j})$.  The Hankel block $H_{n}$ consumes
$m_{0},\dots,m_{2n-2}$, so $n=N=(S+1)/2$ (for odd $S$) is the largest
Hankel block determined entirely by free data: \eqref{eq:lstar} says
the form stays positive definite \emph{through the whole free range
and no further is asked}.  Heuristically, the atoms sit at $S\approx
L/2$ consecutive points of the image of an equispaced $L$-point grid
on the unit circle, and a polynomial of degree $n$ in $x=\cos\theta$
is an even trigonometric polynomial of degree $n$; degrees up to
$(S+1)/2\approx L/4$ are therefore precisely the sub-Nyquist regime
in which $p^{2}$ (of trigonometric degree $\le 2n \approx L/2$) is
just barely unable to localize on individual grid atoms.  Two hard
counting facts frame the problem from above: with $p$ the node
polynomial of $\mu$ one gets $\int p^{2}\dd(\mu-\nu)=-\int
p^{2}\dd\nu<0$, so positive definiteness \emph{must} fail by degree
$S_{+}$; and the number of negative pivots over the full recurrence
is exactly $S_{-}$ (Jacobi/Sylvester).  The observed failure sits at
half-filling, far below the pigeonhole ceiling $S_{+}$ --- counting
alone cannot decide the problem, the content is metric.

\section{Context}\label{sec:context}

The inequality arises in a research program on finite-window versions
of Weil positivity: the lag sequence $c_{i}=A(i\Delta)+P(i\Delta)$ is
the Weil explicit-formula pairing evaluated on triangular test
functions, so $f$ in \eqref{eq:symbol} is an explicit-formula
spectral symbol, and \eqref{eq:lstar} is equivalent to the positive
definiteness of the associated finite quadratic form through the
maximal moment-free range.  The reduction of that program's central
open question to precisely the inequality \eqref{eq:lstar}, together
with the measurement record quoted in \S\ref{sec:numerics}, is
documented in the program paper shipped alongside this note
(\texttt{rh/paper/rh\_program.pdf} in the same repository).  We state
explicitly: nothing here claims progress on the Riemann hypothesis or
on Weil positivity; \eqref{eq:lstar} is an open, self-contained,
finite-dimensional problem in the approximation theory of signed
discrete measures, and both a proof and a counterexample would be
valuable.

\section{The frozen state (August 2026): one object}\label{sec:frozen}

This section was added on August 27, 2026, when the research
programme paused internal work on Problem~\ref{prob:main} and froze
it as a specialist problem.  It compresses the sealed measurement
campaign that followed the update of \S\ref{sec:numerics} (internal
record rounds r342--r354; every number below is quoted from a sealed,
machine-checked record on finite instances).  Nothing in this section
is a theorem about the family, and nothing claims the inequality:
the content is a \emph{reduction} of the open problem to one measured
object and three questions.

\subsection{The extended census and the frame limit}
\label{sec:frozen-census}

Under sealed extension rules (the recipe of \S\ref{sec:construction}
verbatim, the window cap $N\le 900$ lifted, anchors kept under the
domain cap $z^{2}\le 4\cdot10^{5}$) the family was extended in five
tranches to $85$ instances with critical degree $N$ up to $7942$.
On \emph{every} instance the margin $1-\lambda_{\max}(\E_{N})$ is
strictly positive (deepest tranche, $N=6532..7942$:
$1.1\cdot10^{-10}$ to $1.7\cdot10^{-9}$; every near-zero sign
decision in the campaign was arbitrated at $30$--$45$ digits).  Two
frame facts matter for what follows.  (i)~The admissible pool of the
construction is \emph{exhausted} at $N=7942=10^{3.90}$: the census
cannot be extended inside this document's frame.  (ii)~On the deep
tranches the margin decay \emph{flattens} against the power-law fit
of the first $57$ instances --- the deep rows sit uniformly on the
high side of the extrapolation (a sealed two-parameter
running-exponent model puts the census-grade limit exponent near
$1.65$ where the $57$-instance fit gives $3.33$; census only, no law
is claimed --- but the drift is in the direction friendly to the
inequality).

\subsection{The pair reduction and the typed chain}
\label{sec:frozen-chain}

Fix an instance and let $y_{1},y_{2}$ be the two atoms of $\nu$
nearest to $x=1$ (the shallow edge; on the $z=16$ instance these are
the atoms at $j=2,4$).  Write
\[
d_{k}=v_{k}K_{N}(y_{k},y_{k})=(\E_{N})_{kk},\qquad
c=\sqrt{v_{1}v_{2}}\,K_{N}(y_{1},y_{2})=(\E_{N})_{12},
\]
\[
p=1-d_{1},\qquad q=1-d_{2},\qquad
r=1-\frac{c^{2}}{pq}\quad(\text{the pair reserve}).
\]
The extremal direction of $\E_{N}$ concentrates on this pair
uniformly along the family (pair mass $0.935$--$0.999$,
participation ratio $1.61$--$1.94$ on all tested instances), and at
the binding point the inequality is exactly the determinant
condition $c^{2}<pq$ --- a \emph{backward} Cauchy--Schwarz statement
about the $\nu$-dressed kernel at two hard-edge atoms.  The chain of
measured structure above this pair is now fully typed:

\begin{enumerate}
\item \textbf{Weights: closed-form dictionary.}  The pair weights
$v_{1},v_{2}$ are predicted in closed form from the construction ---
archimedean layer through the digamma dictionary
$F_{A}(\xi)=-\log\pi+\operatorname{Re}\psi(\tfrac14+i\xi/2)$ (with an
explicit truncation tail), prime layer through an exact two-cell
tent formula --- at median relative error below $10^{-4}$ (maximum
$9\cdot10^{-4}$) across the family.
\item \textbf{Deficits: candidate laws.}  Along the family the bare
scalars decay as $N^{-a}$ with robust-fit exponents
$a_{p}=0.754$, $a_{q}=0.645$, $a_{c}=0.697$.  The first-hit rational
candidates $a_{p}=\tfrac34$, $a_{q}=\tfrac23$ pass every sealed
cleanliness clause, but a hit is not an identification: separating
$\tfrac34$ from $2\times0.38$ needs a family reaching
$N\sim10^{31.5}$, $\tfrac23$ from $1-0.38$ needs $10^{8.5}$, and the
nearest ambiguity of the chain needs $10^{5.4}$ --- against the pool
end $10^{3.90}$.  The candidates are \emph{unresolvable inside the
frame}.
\item \textbf{Pinning: theorem grade.}  Split
$\E_{N}=\bigl[\begin{smallmatrix}A&C\\C^{\mathsf T}&D\end{smallmatrix}\bigr]$
(pair vs.\ rest) and dress the pair block:
$A'=A+C(I-D)^{-1}C^{\mathsf T}$ with entries $d_{1}',d_{2}',c'$,
deficits $p',q'$, reserve $r'$, and pair margin
$m_{2}'=1-\lambda_{\max}(A')$.  Then
$m_{2}'\,(p'+q'-m_{2}')=p'q'-c'^{2}=p'q'\,r'$ exactly (both dressed
eigenvalue deficits are roots of $t^{2}-(p'+q')t+(p'q'-c'^{2})$); in
the top-2--truncated resolvent model $m_{2}'$ \emph{equals} the full
margin identically and $p'q'-c'^{2}$ equals
$\mathrm{margin}^{2}\cdot g_{21}/\Delta t^{2}$, exact in rational
arithmetic; live, $|m_{2}'/\mathrm{margin}-1|\le0.0605$ on all $75$
fully dressed instances.  Consequently every dressed scalar decays
at the margin rate times an $O(1)$ geometry factor: writing
$\delta_{x}$ for the extra decay of $x'/x$, the identity forces
$\delta_{x}=\alpha-a_{x}$ (matched live at $0.033$), where
$\alpha=3.332$ is the measured margin exponent.
\item \textbf{Exponent bookkeeping closes with one census member.}
$\alpha=a_{p}+a_{q}+\rho_{r}-a_{p+q}
=0.754+0.645+2.624-0.690=3.333$ against the measured $3.332$
(deviation $0.001$); with the rational candidates,
$\tfrac34+\tfrac23+\rho_{r}-a_{p+q}=3.352$ (deviation $0.019$).
Here $\rho_{r}=2.624$, the decay exponent of the reserve $r$, is the
single member of the chain that is neither dictionary nor candidate
nor theorem: it is a census fit over a visibly curved family
(halves curvature $-0.767$; no clean power law), and \S\ref
{sec:frozen-object} identifies what it actually is.
\item \textbf{Computability closure.}  The composition
$c_{\mathrm{pred}}
=\sqrt{v_{\mathrm{pred},1}v_{\mathrm{pred},2}}\;K_{12}^{\mathrm{cd}}$
--- closed-form dictionary weights times the cross kernel evaluated
by the Christoffel--Darboux recursion from the source measure, with
no measured kernel input --- reconstructs $c$ on all $85$ instances
at maximum relative deviation $5.5\cdot10^{-4}$ and predicts the
fine structure $\varphi$ of \S\ref{sec:frozen-object} pointwise at
correlation $1.0000000$ and residual rms ratio $0.0004$.  This is a
\emph{computability} statement, not a closed form and not a
derivation.
\end{enumerate}

\subsection{The one object: a shared wander that cancels}
\label{sec:frozen-object}

Set $A_{\mathrm{lead}}=\tfrac{17}{12}=\tfrac34+\tfrac23$ (the sealed
common leading exponent, the sum of the first-hit candidates) and
define, across the family, the two fine structures
\[
\varphi_{D}=\log(pq)+A_{\mathrm{lead}}\ln N,
\qquad
\varphi_{K}=\log(c^{2})+A_{\mathrm{lead}}\ln N .
\]
The frozen finding (measured on the $81$-instance census, reproduced
exactly on the $85$):
\begin{itemize}
\item $\varphi_{D}$ and $\varphi_{K}$ are \emph{one shared wander}:
correlation $0.999998$, each with rms $0.88$ nats, while their
difference has rms $0.0017$ nats --- a factor $\sim500$ smaller.
\item The difference is not an approximation --- it is pointwise the
log reserve: $\varphi_{D}-\varphi_{K}=-\log(c^{2}/pq)=-\log(1-r)$,
and its decay refits $\rho_{r}$ at $|{\Delta}\text{slope}|=0.0002$.
\emph{The one unexplained exponent of the chain, $\rho_{r}=2.624$,
is exactly the decay of the difference of two near-identical
wanders.}
\item The wander has a measured anatomy.  Via
$c=\sqrt{v_{1}v_{2}}\,K_{N}(y_{1},y_{2})$ and the
Christoffel--Darboux endpoint formula, $\log c^{2}$ factorizes
pointwise into four source blocks,
\[
\log c^{2}=W+B_{2}+P_{2}+G_{M},
\]
with weights $W=\log(v_{1}v_{2})$, recurrence prefactor
$B_{2}=2\log b_{N-1}$, CD polynomial
$P_{2}=2\log|p_{N}(y_{1})p_{N-1}(y_{2})-p_{N-1}(y_{1})p_{N}(y_{2})|$,
and gap $G_{M}=-2\log|y_{1}-y_{2}|$.  No single block carries
$\varphi$: the pair-geometry block ($P_{2}$, equivalently the
composite $P_{2}+G_{M}$) matches its \emph{shape} at correlation
$+0.887$ but with $2.4\times$ the amplitude, and the weight block
$W$ is \emph{anti}-correlated at $-0.72$ and cancels the excess.
$\varphi$ is a two-block interference --- a large pair-geometry
wander almost cancelled by the dictionary-weight wander.
\item $\varphi$ is pair-local: it is not a function of $N$ alone
(instances of equal depth disagree at the full wander scale,
adjacency ratio $0.911$), and neither the local gap class nor the
anchor size carries it (residual ratios $0.81$/$0.91$).
\end{itemize}
The frozen question is therefore a single one: \textbf{why do the
pair-geometry wander and the dictionary-weight wander cancel to a
residual five hundred times smaller than either --- equivalently,
why does $c$ track $\sqrt{pq}$ at two hard-edge atoms to this
precision, with the cancellation decaying at the specific rate
$\rho_{r}=2.624$?}

\subsection{Three questions for specialists (final form)}
\label{sec:frozen-questions}

\begin{enumerate}
\item[(q1)] \textbf{The backward Cauchy--Schwarz fine structure ==
the cancellation ratio.}  For the $\nu$-dressed
Christoffel--Darboux kernel of $\mu$ at the two shallow-edge atoms,
explain the near-saturation $c^{2}/pq\to1$ at rate $N^{-2.624}$ ---
in the sharpest measured form: is the two-block cancellation ratio
of \S\ref{sec:frozen-object} (pair geometry against dictionary
weights, $2.4\times$ amplitude, correlation $+0.887$ against
$-0.72$) a law?
\item[(q2)] \textbf{The sub-classical Christoffel growth.}  The
within-window growth of $K_{n}(y,y)$ at the shallow-edge atoms has
halves exponent $\approx0.38$ (census $0.336$--$0.608$) ---
sub-classical for a hard edge --- and the across-family deficit
exponents hit the rational candidates $\tfrac34,\tfrac23$ and the
$0.38$-family ($2\times0.38$, $1-0.38$) simultaneously.  Deciding
between them needs families reaching $N\sim10^{5.4}$ (nearest
ambiguity) up to $10^{31.5}$, while this construction ends at
$10^{3.90}$.  What is the true growth family, and why?
\item[(q3)] \textbf{The resolution paradox.}  At the critical degree
the kernel resolves the pair, and the bare reserve saturates
($r\to0$ as $N^{-2.6}$); yet the \emph{dressed} reserve $r'$ stays
$O(1)$ (span $0.138$--$0.536$ over $75$ instances, certified flat by
a curvature-honest protocol; likewise the spectral gap ratio
$g_{21}=(1-\lambda_{2})/(1-\lambda_{1})$, span $3.4$--$19.7$).  The
dressing works by \emph{signed cancellation}, $\Delta c\approx-c$:
every absolute-value bound chain (triangle, Cauchy--Schwarz,
operator norm) certifies none of the $75$ instances, overshooting by
factors $73$ up to $2.4\cdot10^{11}$.  Why does the dressing pin the
dressed block at $O(1)$ --- what controls the signed cancellation?
\end{enumerate}

\subsection{The honest frame}\label{sec:frozen-status}

Everything above is a census on $85$ finite instances with critical
degrees up to $N=7942$; no statement is asymptotic.  All measured
margins are strictly positive --- no counterexample --- and no proof
mechanism is known; the inequality itself remains fully open, for
the unbounded family in both directions.  The candidate exponents
are unresolvable inside this frame and the admissible pool is
exhausted, which is precisely why the problem is frozen here rather
than measured further: additional internal numerics on this family
cannot decide anything that is still open.  The record rounds cited
above (r342--r354) are sealed, hash-pinned probe records in the
research repository that ships this document; the companion script
of \S\ref{sec:script} is unchanged and remains the machine check of
the construction.  This section, like the rest of the document,
claims nothing: not the inequality, and nothing about the Riemann
hypothesis.

\subsection{Placement in the literature, and addressees}
\label{sec:frozen-placement}

This closing subsection (added on the evening of August 27, 2026)
translates the frozen object into the standard language of the
communities whose toolkits contain it.  It adds no new measurement
and no claim; the classification below is census plus one clearly
marked hypothesis.

\subsubsection*{Dictionary: the multi-point Uvarov transform}

The signed measure $\mu-\nu$ is, verbatim, a multi-point
(generalized) \emph{Uvarov transform} of the positive measure
$\mu$: the transform subtracts the point masses $v_{k}$ at the
nodes $y_{k}$.  Quasi-definiteness of such transforms is a
classical subject.  Kwon et al.\ (\emph{Quasi-definiteness of
generalized Uvarov transforms of moment functionals}, J.~Appl.\
Math.\ 1 (2001)) characterize quasi-definiteness of the
transformed functional by the nonvanishing of determinants
$\det D_{n}\neq0$ built from the Christoffel--Darboux kernel
values $K_{n}(y_{k},y_{l})$ at the mass points;
Akta\c{s}--Area--P\'erez (Comput.\ Appl.\ Math.\ 41 (2022),
Thm~4) state the same criterion in matrix form: the transform is
quasi-definite at degree $n$ iff $I+\Lambda K_{n}$ is
nonsingular, with $\Lambda$ the diagonal matrix of the added
masses and $K_{n}$ the kernel matrix at the mass points.  With
$\Lambda=-\operatorname{diag}(v_{k})$ the matrix $I+\Lambda
K_{n}$ is similar to $I-\E_{n}$ of \S\ref{sec:equiv}(b)
(conjugate by $\operatorname{diag}(\sqrt{v_{k}})$).  The objects
of this document are therefore exactly the standard ones:
$\E_{N}$ is the multi-point Uvarov kernel matrix of the
modification at the $\nu$ atoms; the free-window
quasi-definiteness of \S\ref{sec:equiv}(a) is the standard
quasi-definiteness question for this transform, asked
simultaneously at all degrees $n\le N$; and the pair reserve
$r=1-c^{2}/(pq)$ of \S\ref{sec:frozen-chain} is a normalization
of the $2\times2$ Uvarov determinant at the two binding atoms.
What is \emph{not} standard: the masses and nodes here are
prime-generated rather than free parameters; the inequality asks
the definite (spectral, $\lambda_{\max}<1$) refinement of
quasi-definiteness, uniformly through the moment-free window;
and the frozen content is the measured \emph{rate} at which the
$2\times2$ determinant closes ($\rho_{r}=2.624$).

\subsubsection*{Universality class: the discrete hard edge
(census, one marked hypothesis)}

Measured: the weight-free kernel correlation at the pair,
$\rho_{K}=K_{N}(y_{1},y_{2})^{2}/\bigl(K_{N}(y_{1},y_{1})\,
K_{N}(y_{2},y_{2})\bigr)$, \emph{decays} along the family with
fitted exponent $1.42$ (record round r352; exactly twice the
fitted exponent $0.711$ of
$\kappa=K_{12}/\sqrt{K_{11}K_{22}}$, the fit homogeneity
exact).  The two atoms sit at hard-edge-rescaled positions
\[
a_{j}=2N^{2}\bigl(1-y_{j}\bigr),
\qquad
a_{j}\;\longrightarrow\;\frac{\pi^{2}j^{2}}{4}
\quad(j\in\{2,4\}\ \text{fixed},\ N\to\infty),
\]
asymptotically family-constant by the grid recipe
($y_{j}=\cos(2\pi j/L)$, $L=4N-2$); on the running instance
$a_{2}=9.92$ and $a_{4}=39.69$, against the limits
$\pi^{2}=9.87$ and $4\pi^{2}=39.48$.  If $\mu$ belonged to the
absolutely-continuous hard-edge universality class (Lubinsky,
Int.\ Math.\ Res.\ Papers, 2008), the $2N^{2}$-rescaled kernel
would converge to the Bessel kernel, and $\rho_{K}$ at these
fixed rescaled positions would tend to a \emph{positive
constant}.  It decays instead: the family is measurably outside
the AC hard-edge class --- consistent with the construction,
since $\mu$ is purely atomic and the AC hypotheses do not apply.
The edge lives in the \emph{discrete} class.  Within that class
the natural regime is the saturation regime of
Baik--Kriecherbauer--McLaughlin--Miller (\emph{Discrete
Orthogonal Polynomials: Asymptotics and Applications}, Ann.\ of
Math.\ Studies 164, Princeton, 2007): the kernel is built from
the orthonormal polynomials of the discrete measure $\mu$ used
at degree $N\approx0.70\,S_{+}$ (on the running instance
$184/263=0.6996$), and in that theory the edge parametrix
adjacent to a saturated band is built out of Euler Gamma
functions (Thm~2.7 and Remark~2.8 there) --- noteworthy here
because the closed-form weight dictionary of
\S\ref{sec:frozen-chain} is the digamma $\psi=(\log\Gamma)'$.
Status of this placement: the decorrelation and the rescaled
positions are \emph{measured}; the saturation-regime adjacency
is a \emph{hypothesis} (an internal round in progress probes the
duality/regime side), and no asymptotic statement is claimed.

\subsubsection*{Who might recognize this}

Four communities own pieces of this object.  The discrete
orthogonal-polynomial asymptotics school (the
Baik--Kriecherbauer--McLaughlin--Miller line): Riemann--Hilbert
asymptotics of discrete orthogonal polynomials, hard edges next
to saturated bands --- question (q2) and the class placement
above live there.  The Christoffel-function and universality
community (the Lubinsky--Levin--Totik line): growth of
$K_{n}(y,y)$ at edge atoms of purely discrete measures --- (q2)
is a Christoffel-function question, and the measured
decorrelation is a universality-class datum.  The
Uvarov/spectral-transformation circle (the Marcell\'an school):
$\E_{N}$ and the pair determinant are their standard objects
(dictionary above); (q1) and (q3) are statements about
multi-point Uvarov determinants and their dressed $2\times2$
blocks.  The Weil-positivity environment (the Connes--Consani
line): truncated and finite-rank forms of Weil positivity; the
finite Guinand--Weil dictionary of arXiv:2607.02828 (2026) is a
related finite tool through which the cancellation of
\S\ref{sec:frozen-object} could be re-expressed.  These are
fields, not assignments: the paragraph names the schools whose
standard toolkits contain the objects of this document, and it
claims nothing on their behalf.

\section{Verification script, contact, citation}\label{sec:script}

The companion script \texttt{verify\_lstar\_instance.py} (same
directory; \texttt{numpy}/\texttt{mpmath} only) rebuilds
$(\mu,\nu)$ for $z=16$ from the formulas of
\S\ref{sec:construction} alone, cross-checks every atom position and
weight against the research code that produced the measurement record
(agreement to $4\cdot10^{-16}$ relative), reproduces
$\lambda_{\max}(\E_{184})=0.99983248$ and
$\lambda_{\max}(\E_{185})=1.00003660$ with a $60$-digit ward, and
verifies \eqref{eq:lstar} directly on monomials and on $500$ random
polynomials of degree $<184$; \texttt{--ladder} sweeps all $42$
instances (the margin table of \S\ref{sec:numerics} is its output).
It ends \texttt{LSTAR INSTANCE VERIFIED}.

\medskip\noindent
Questions, proof sketches and counterexamples are equally welcome.
Suggested citation: S.~Hamann, \emph{A strict polynomial
subordination inequality for two explicit prime-generated discrete
measures (open problem)}, TFPT research repository,
\texttt{rh/problem/}, August 2026.

\end{document}
